\section{Система непересекающихся множеств (DSU)}
Поддерживает объединение компонент и проверку принадлежности одной компоненте; объединение возвращает признак изменения разбиения.
Вход --- векторы \texttt{parent} и \texttt{size} размера $n$, исходно $parent[i]=i$, $size[i]=1$, вершины $0\ldots n-1$; требуется $0\le n<\texttt{INT\_MAX}$ и корректные индексы вершин.
Поиск корня и объединение занимают амортизированно $O(\alpha(n))$ времени, память $O(n)$.
\begin{lstlisting}[name=data_dsu]
int find_root(vector<int>& parent, int vertex) {
    if (parent[vertex] == vertex) return vertex;
    return parent[vertex] = find_root(parent, parent[vertex]);
}

bool unite(vector<int>& parent, vector<int>& size, int first, int second) {
    first = find_root(parent, first);
    second = find_root(parent, second);
    if (first == second) return false;
    if (size[first] < size[second]) swap(first, second);
    parent[second] = first;
    size[first] += size[second];
    return true;
}
\end{lstlisting}
% \subsection*{Для чего нужен}
% Поддерживает объединение компонент и проверку принадлежности одной компоненте.
% Вершины нумеруются от $0$ до $n-1$, векторы \texttt{parent} и \texttt{size}
% имеют размер $n$; исходно $parent[i]=i$, $size[i]=1$.
% Требуется $0\le n<\texttt{INT\_MAX}$.
% Амортизированное время поиска корня и объединения ---
% $O(\alpha(n))$, память $O(n)$. Объединение возвращает признак
% изменения разбиения.
% Код адаптирован из \url{https://ru.algorithmica.org/cs/set-structures/dsu/}:
% массивы передаются параметрами, добавлены проверка совпадения корней
% и возвращаемый признак объединения.
% \subsection*{Реализация}
%
% \subsection*{Возможные улучшения и дополнения}
% Сжатие путей и объединение по размеру дают почти константную амортизацию.
% Можно хранить дополнительные данные компоненты в корне и объединять их при
% смене корня. Для отката операций применяют DSU с историей изменений без
% сжатия путей. Для офлайн-задач полезен DSU на отрезке времени.

\section{Дерево Фенвика}
Поддерживает прибавление к позиции и суммы на полуинтервалах: префикс относится к $[0,right)$, запрос суммы --- к $[left,right)$.
Вход --- изначально нулевой вектор \texttt{tree} размера $n+1$; требуется $0\le n<\texttt{INT\_MAX}$, $0\le left\le right\le n$, при обновлении $0\le index<n$.
Суммы и промежуточные значения должны помещаться в \texttt{ll}; каждая операция занимает $O(\log(n+1))$ времени и $O(1)$ дополнительной памяти.
\begin{lstlisting}[name=data_fenwick]
void add(vector<ll>& tree, int index, ll delta) {
    int length = (int)tree.size() - 1;
    for (++index; index <= length;) {
        tree[index] += delta;
        int step = index & -index;
        if (step > length - index) break;
        index += step;
    }
}

ll prefix(const vector<ll>& tree, int right) {
    ll answer = 0;
    for (; right > 0; right -= right & -right) answer += tree[right];
    return answer;
}

ll range(const vector<ll>& tree, int left, int right) {
    return prefix(tree, right) - prefix(tree, left);
}
\end{lstlisting}
% \subsection*{Для чего нужен}
% Поддерживает прибавление к одной позиции и суммы на полуинтервалах.
% Внешний вектор \texttt{tree} имеет размер $n+1$ и изначально заполнен нулями;
% индексы элементов --- $0\ldots n-1$. Префиксная сумма относится
% к $[0,right)$, сумма отрезка --- к $[left,right)$.
% Требуется $0\le left\le right\le n$, для обновления $0\le index<n$;
% $n<\texttt{INT\_MAX}$. Суммы и промежуточные значения помещаются в \texttt{ll}.
% Каждая операция занимает $O(\log(n+1))$ времени и $O(1)$ дополнительной памяти.
% Код адаптирован из \url{https://ru.algorithmica.org/cs/range-queries/fenwick/}:
% внешний массив передаётся параметром, индексы элементов начинаются с нуля,
% границы полуоткрытые, суммы имеют тип \texttt{ll}.
% \subsection*{Реализация}
%
% \subsection*{Возможные улучшения и дополнения}
% Для значения в одной позиции и прибавления на отрезке достаточно одного
% дерева Фенвика над разностями. Два дерева позволяют получать суммы при
% прибавлении на отрезке. Для минимума возможны специальные варианты только при ограниченных
% обновлениях; для произвольных изменений удобнее дерево отрезков.

\section{Дерево отрезков с отложенным прибавлением}
Обрабатывает массив значений и запросы $(left,right,value,type)$: тип $0$ прибавляет \texttt{value} на $[left,right)$, тип $1$ возвращает сумму, игнорируя поле значения; ответ содержит суммы в порядке запросов.
Требуется $0\le n<\texttt{INT\_MAX}/4$, $0\le left\le right\le n$; пустые массивы и полуинтервалы допустимы, суммы, метки и промежуточные произведения должны помещаться в \texttt{ll}.
Построение занимает $O(n)$, запрос --- $O(\log(n+1))$, общая память $O(n+q)$.
\begin{lstlisting}[name=data_segment]
vector<ll> segmentTree(const vector<ll>& values,
        const vector<tuple<int, int, ll, int>>& queries) {
    int n = values.size();
    vector<ll> tree(4 * max(1, n)), lazy(tree.size()), answer;
    auto build = [&](auto&& self, int node, int left, int right) -> void {
        if (right - left == 1) { tree[node] = values[left]; return; }
        int middle = (left + right) / 2;
        self(self, node * 2, left, middle);
        self(self, node * 2 + 1, middle, right);
        tree[node] = tree[node * 2] + tree[node * 2 + 1];
    };
    auto update = [&](auto&& self, int node, int left, int right,
                      int queryLeft, int queryRight, ll value) -> void {
        if (queryRight <= left || right <= queryLeft) return;
        if (queryLeft <= left && right <= queryRight) {
            tree[node] += value * (right - left); lazy[node] += value; return;
        }
        int middle = (left + right) / 2;
        self(self, node * 2, left, middle, queryLeft, queryRight, value);
        self(self, node * 2 + 1, middle, right, queryLeft, queryRight, value);
        tree[node] = tree[node * 2] + tree[node * 2 + 1]
                   + lazy[node] * (right - left);
    };
    auto query = [&](auto&& self, int node, int left, int right,
                     int queryLeft, int queryRight, ll carry) -> ll {
        if (queryRight <= left || right <= queryLeft) return 0;
        if (queryLeft <= left && right <= queryRight)
            return tree[node] + carry * (right - left);
        int middle = (left + right) / 2;
        carry += lazy[node];
        return self(self, node * 2, left, middle, queryLeft, queryRight, carry)
             + self(self, node * 2 + 1, middle, right, queryLeft, queryRight, carry);
    };
    if (n) build(build, 1, 0, n);
    for (auto [left, right, value, type] : queries)
        if (type == 0) { if (n) update(update, 1, 0, n, left, right, value); }
        else answer.push_back(n ? query(query, 1, 0, n, left, right, 0) : 0);
    return answer;
}
\end{lstlisting}
% \subsection*{Для чего нужен}
% Пакетно обрабатывает прибавления на полуинтервале и запросы суммы.
% Запрос хранит левую и правую границы, значение и тип в этом порядке:
% при типе $0$ значение прибавляется на $[left,right)$, при типе $1$
% вычисляется сумма, а поле значения игнорируется. Ответ содержит суммы
% в порядке запросов. Требуется $0\le left\le right\le n$;
% пустые массивы и полуинтервалы допустимы.
% Построение занимает $O(n)$, запрос --- $O(\log(n+1))$,
% общая память $O(n+q)$. Требуется $n<\texttt{INT\_MAX}/4$;
% суммы, метки и промежуточные произведения помещаются в \texttt{ll}.
% Тематический источник: \url{https://ru.algorithmica.org/cs/segment-tree/lazy-propagation/}.
% Код ниже самостоятельный: в статье используется присваивание на отрезке
% и дерево на указателях, здесь --- прибавление и массивы без проталкивания.
% \subsection*{Реализация}
%
% \subsection*{Возможные улучшения и дополнения}
% При необходимости можно вернуть минимум или максимум, заменив агрегат и
% формулу применения отложенного значения. Для онлайн-запросов вынесите
% построение и операции в функции с внешними \texttt{tree} и \texttt{lazy}.
% Комбинацию прибавления и присваивания нужно хранить двумя видами меток и
% применять их в правильном порядке.

